% SPDX-License-Identifier: MIT
\chapter{Fehlersuche und bewusste Grenzen}\label{chap:trouble}
\ChapterSummary{Dieses Kapitel liefert eine symptomorientierte
Diagnose für die häufigsten Probleme bei Build,
Terminalbetrieb, Layout, Eingabe und Modellbindung.
Es zeigt ausserdem, welche Eigenschaften derzeit
noch keine garantierte Produktfunktion sind.}

\section{CMake findet das Toolkit nicht}
\textbf{Symptom:} \texttt{find\_package} meldet
kein \texttt{SASDUIToolkitConfig.cmake}.

\textbf{Prüfen:} Wurde die Library wirklich
installiert? Zeigt \filepath{CMAKE_PREFIX_PATH}
auf den \emph{Installationspraefix}, nicht
auf eine beliebige Quelldatei?
Fordert der Konsument eine zum
installierten Paket passende Version?
Sind Quell-Build und bereits
installierte alte Versionen verwechselt?
Ein sauberer Consumer-Build in einem
neuen Verzeichnis beseitigt viele
Caches mit falschen Fundstellen.

\section{Unter Visual Studio wird nichts gestartet}
\textbf{Symptom:} Sie haben kompiliert,
finden aber das Beispiel nicht.

\textbf{Prüfen:} Der Visual-Studio-
Generator nutzt Konfigurationsordner
wie \filepath{Debug} und
\filepath{Release}. Wurde die
richtige Architektur gewählt?
Passen \filepath{--config Debug}
bei Build und \filepath{-C Debug}
bei CTest zusammen?
Bei einem abweichenden Generator
kann der Pfad anders aussehen.

\section{Terminal bleibt leer oder zeigt Zeichenmuell}
\textbf{Symptom:} Beim Demo-Start
erscheint kein sinnvoller Bildschirm.

\textbf{Prüfen:} Laeuft die Anwendung
in einem interaktiven, ANSI/VT-fähigen
Terminal? Werden Ein-/Ausgabe eventuell
in eine Datei oder fremde Prozesskonsole
umgelenkt? Ist das Terminal sehr klein?
Unterstützt das Terminal die
benötigten Escape-Sequenzen?
Testen Sie die Referenz-Demo, bevor
Sie den eigenen Host untersuchen.
Der Core allein sendet keine Frames;
die Kombination aus PresentationSink,
Frame-Encoder und Backend ist notwendig.

\section{Button reagiert nicht}
\textbf{Prüfen:} Ist er sichtbar,
aktiviert und für die gewuenschte
Bedienung fokussierbar?
Erhaelt der Host die KeyEvents
und leitet er sie an das korrekte
Fokusziel weiter?
Ist ein Command gebunden und
möglicherweise deaktiviert?
Hat ein geöffnetes Menue oder
Popup den Eingabefokus?
Bei Mausaktionen ist zusaetzlich
Pointer-Routing/Capture zu prüfen.

\section{Nach Textaenderung sieht man nichts}
\textbf{Prüfen:} Verändert
\texttt{setText()} wirklich den
Wert? Fordert die Host-Schleife
nach Invalidierung einen
Praesentations-Refresh an?
Wurde die Größe bei längerem
Text erneut gemessen und
arrangiert?
Handelt es sich um ein
verdecktes oder geclipptes Widget?
Die Kernobjekte sind kein
automatisches Desktop-Fenstersystem.

\section{Tab-Reihenfolge oder Eingabe stimmt nicht}
\textbf{Prüfen:} Wurde
\term{FocusTraversal} mit
dem aktuellen sichtbaren Baum
benutzt? Können die betreffenden
Widgets Fokus erhalten?
Sind TextInputEvents von
KeyEvents getrennt?
Leitet ein globaler Shortcut
irrtuemlich dieselbe Taste weiter?
Wurde ein altes Popup oder
Pointer-Capture nicht geschlossen?

\section{Das Model zeigt alte oder falsche Zeilen}
\textbf{Prüfen:} Wurde ein
mitgeliefertes StringModel über
seine Mutationsmethoden
aktualisiert? Meldet ein
eigenes Model seine Änderungen
und Revisionen korrekt?
Sind bei der Table alle Zeilen
rechteckig? Wurde
die Auswahl nach strukturellen
Änderungen neu validiert?
Ist der Viewport noch passend?
Wurden die Lebensdauern von Model,
SelectionModel und View bedacht?

\section{Unterschiedliche Darstellung unter SDL3}
\textbf{Prüfen:} Ist der
optionale Adapter wirklich aktiviert?
Sind SDL3 und SDL\_ttf in der
geforderten Version vorhanden?
Ist ein gültiger Fontpfad
für entsprechende Tests gesetzt?
Arbeitet die normale Rendered-
Schicht im Headless-Test korrekt?
SDL3 ist gegenwärtig
experimentell und nicht mit
voller Terminal-/Native-Paritaet
gleichzusetzen.

\section{Was noch nicht zum Vertrag gehört}
Das Toolkit hat derzeit keinen
abgeschlossenen visuellen RAD-Designer,
keinen garantierten universellen
\texttt{Application::run()}-Host,
keine fertigen nativen Win32-/GTK-/AppKit-
Peers und keine Zusage, dass jedes
Core-Widget auf allen Backends
identisch bedienbar ist.
Auch grosse Daten-Views und
komplexe Popup-Interaktion sollten
im Zielbackend stets mit den
tatsächlichen Demos und Tests
abgeglichen werden.

\section{Wie Sie einen hilfreichen Fehlerbericht schreiben}
Ein guter Bericht nennt den exakten
Commit/Tag, Compiler, Plattform,
Terminal oder SDL-Konfiguration,
CMake-Flags und die kleinste
Reproduktionsfolge. Fuegen Sie
einen kurzen, vollständigen
Codeausschnitt und die
\emph{erste} entscheidende
Fehlerausgabe bei.
Beschreiben Sie erwartet und
beobachtet getrennt.
Vor einem Pull Request sollten
die passenden CTest-Suiten laufen.
Technische Diskussionen finden
unter \url{https://github.com/Robin-Goerlach/SASD-UI-Toolkit/issues} statt.

\PracticalTip{Bei einer unklaren Stoerung
zuerst das mitgelieferte Beispiel
auf genau demselben System starten,
dann Ihr Programm schrittweise
auf denselben Komponentenpfad
reduzieren. Der Unterschied zeigt
oft bereits die fehlende Host-
oder Lifecycle-Verbindung.}
